\section{Empirical Evaluations}
\label{sec:evaluation}

We evaluate DeXposure-FM on the DeXposure dataset, focusing on three tasks: (i) \textit{multi-step forecasting} of edge-level exposures and network-level statistics, (ii) \textit{shock analysis} during major stress events (Terra/Luna, FTX), and (iii) \textit{imputation} of missing edges.

\subsection{Dataset and Experimental Setup}
\label{subsec:dataset}

\paragraph{Data source.}
We use the DeXposure dataset constructed from DeFiLlama API data, spanning March 2020 to August 2025 (283 weekly snapshots). Each snapshot contains protocol nodes with TVL measurements and directed edges representing credit exposures between protocols.

\paragraph{Data statistics.}
Table~\ref{tab:data-stats} summarizes the dataset characteristics. The network exhibits high edge persistence (mean overlap ratio 98.5\%), reflecting the stable nature of DeFi protocol interconnections.

\begin{table}[t]
\centering
\small
\begin{tabular}{lr}
\hline
\textbf{Statistic} & \textbf{Value} \\
\hline
Total snapshots & 283 weeks \\
Time range & 2020-03-23 $\sim$ 2025-08-18 \\
Mean nodes per week & 5,676 \\
Mean edges per week & 30,424 \\
Mean edge overlap ratio & 98.5\% \\
Protocol categories & $\sim$15 classes \\
\hline
\end{tabular}
\caption{DeXposure dataset statistics.}
\label{tab:data-stats}
\end{table}

\paragraph{Temporal split.}
We use expanding window walk-forward validation with a 2025 hold-out test set:
\begin{itemize}
  \item \textbf{Training:} 2020-03-23 $\sim$ 2024-10-07 (238 weeks)
  \item \textbf{Validation:} 2024-10-14 $\sim$ 2024-12-30 (12 weeks)
  \item \textbf{Test:} 2025-01-06 $\sim$ 2025-08-18 (33 weeks, hold-out)
\end{itemize}

\paragraph{Models compared.}
We compare three approaches:
\begin{enumerate}
  \item \textbf{GraphPFN (Fine-tuned):} Pre-trained GraphPFN encoder with end-to-end fine-tuning using differential learning rates.
  \item \textbf{ROLAND:} Temporal GNN baseline with GCN encoder and GRU temporal updates, trained from scratch.
\end{enumerate}

\subsection{Task I: Multi-step Forecasting}
\label{subsec:forecast}

\paragraph{Targets.}
We evaluate forecasts at horizons $h \in \{1, 3, 7\}$ weeks for:
\begin{enumerate}
  \item \textbf{Edge existence:} Binary prediction of whether edge $(p,q)$ exists at $t+h$.
  \item \textbf{Edge weight:} Regression on log-scaled edge weight for existing edges.
  \item \textbf{Node TVL change:} Regression on $\Delta_{p,t+h} = \log(1+\text{TVL}_{t+h}) - \log(1+\text{TVL}_t)$.
\end{enumerate}

\paragraph{Metrics.}
For edge existence, we report AUPRC (area under precision-recall curve) and AUROC. For edge weight and node predictions, we report MAE and RMSE on log-scaled values.

\paragraph{Results.}
Table~\ref{tab:main-results} presents the main forecasting results on the 2025 hold-out test set.

\begin{table}[t]
\centering
\small
\setlength{\tabcolsep}{4pt}
\begin{tabular}{llcccccc}
\hline
& & \multicolumn{2}{c}{\textbf{Edge Exist}} & \multicolumn{2}{c}{\textbf{Edge Weight}} & \multicolumn{2}{c}{\textbf{Node $\Delta$TVL}} \\
\textbf{Model} & $h$ & AUPRC & AUROC & MAE & RMSE & MAE & RMSE \\
\hline
\multirow{3}{*}{GraphPFN (FT)} 
& 1 & \textbf{0.932} & \textbf{0.985} & \textbf{2.62} & \textbf{3.53} & 0.060 & 0.401 \\
& 3 & \textbf{0.934} & \textbf{0.986} & \textbf{2.66} & \textbf{3.56} & 0.121 & 0.605 \\
& 7 & \textbf{0.936} & \textbf{0.986} & \textbf{2.64} & \textbf{3.55} & 0.211 & 0.840 \\
\hline
\multirow{3}{*}{ROLAND}
& 1 & 0.870 & 0.962 & 3.99 & 5.84 & \textbf{0.057} & \textbf{0.403} \\
& 3 & 0.866 & 0.959 & 3.97 & 5.82 & \textbf{0.114} & \textbf{0.608} \\
& 7 & 0.861 & 0.955 & 3.94 & 5.79 & \textbf{0.195} & \textbf{0.847} \\
\hline
\end{tabular}
\caption{Multi-step forecasting results on 2025 hold-out test set. Bold indicates best performance. GraphPFN (FT) = Fine-tuned GraphPFN. Edge weight MAE/RMSE are on log-scaled values.}
\label{tab:main-results}
\end{table}

\paragraph{Key findings.}
\begin{enumerate}
  \item \textbf{Edge existence:} GraphPFN achieves AUPRC $>$0.93 across all horizons, outperforming ROLAND by 6--7 percentage points. The high baseline performance reflects the 98.5\% edge overlap ratio in the data.
  \item \textbf{Edge weight:} GraphPFN reduces MAE by 34\% compared to ROLAND (2.6 vs. 3.9), demonstrating superior weight prediction capability.
  \item \textbf{Node prediction:} ROLAND performs slightly better on node TVL change prediction, possibly due to its explicit temporal modeling via GRU.
  \item \textbf{Horizon stability:} GraphPFN maintains stable performance across horizons (h=1 to h=7), while ROLAND shows mild degradation.
\end{enumerate}

\subsection{Task II: Shock Analysis}
\label{subsec:shock}

We evaluate model behavior during two major DeFi stress events:
\begin{enumerate}
  \item \textbf{Terra/Luna collapse} (May 9, 2022): UST stablecoin de-peg triggering cascading liquidations.
  \item \textbf{FTX collapse} (November 7, 2022): Exchange failure and contagion through FTT token exposures.
\end{enumerate}

For each event, we analyze a 4-week window around the shock date, examining:
\begin{itemize}
  \item Network structure changes (TVL, edge count, concentration metrics)
  \item Model prediction performance during stress periods
  \item Protocol-level impact assessment
\end{itemize}

\subsection{Task III: Imputation}
\label{subsec:imputation}

We evaluate the model's ability to reconstruct missing edges under random masking.

\paragraph{Masking protocol.}
We apply Missing Completely at Random (MCAR) masking at rates $\{10\%, 20\%, 30\%\}$ to the edge set and evaluate reconstruction quality.

\paragraph{Metrics.}
We report AUPRC for edge existence reconstruction and MAE for edge weight reconstruction on masked edges.

\subsection{Ablation: Encoder Comparison}
\label{subsec:ablation}

Table~\ref{tab:ablation} compares GraphPFN against ROLAND to isolate the contribution of the pre-trained foundation model.

\begin{table}[t]
\centering
\small
\begin{tabular}{lcc}
\hline
\textbf{Model} & \textbf{AUPRC} ($h$=1) & \textbf{Weight MAE} ($h$=1) \\
\hline
GraphPFN (Fine-tuned) & \textbf{0.932} & \textbf{2.62} \\
ROLAND (from scratch) & 0.870 & 3.99 \\
\hline
$\Delta$ (improvement) & +7.1\% & $-$34.3\% \\
\hline
\end{tabular}
\caption{Ablation comparing pre-trained GraphPFN vs. ROLAND trained from scratch.}
\label{tab:ablation}
\end{table}

The pre-trained GraphPFN encoder provides substantial gains over training from scratch, validating the transfer learning approach for DeFi exposure networks.

\subsection{Discussion}
\label{subsec:discussion}

\paragraph{High edge overlap and baseline accuracy.}
The 98.5\% mean edge overlap ratio explains the high AUPRC values---most edges persist week-to-week. However, GraphPFN's advantage over ROLAND demonstrates its superior ability to identify the 1.5\% of edges that change.

\paragraph{Practical implications.}
The strong edge weight prediction (34\% lower MAE) is particularly valuable for risk monitoring, as it enables more accurate estimation of exposure magnitudes for macro-prudential analysis.

\paragraph{Limitations.}
Current experiments use h $\in \{1, 3, 7\}$; the h=14 horizon is pending. Future work includes implementing the full 7-loss-component objective and TVL-weighted economic importance sampling as specified in the training design.
